\documentclass[12pt,letterpaper]{scrartcl}
\usepackage{lipsum}
\usepackage[utf8]{inputenc}
\usepackage{amsmath}
\usepackage{amsfonts}
\usepackage{amssymb}
\usepackage{graphicx}
\usepackage[left=3cm,right=2.5cm,top=2.5cm,bottom=2.5cm]{geometry}
\usepackage[]{algorithm2e}
\author{Don cuyi}

%Color
\usepackage{color}
\definecolor{nred}{RGB}{174,49,54}
\definecolor{nblue}{RGB}{86,99,146}
\definecolor{nalgo}{RGB}{188,139,76}
\usepackage{sectsty}
\sectionfont{\color{nred}}
\subsectionfont{\color{nblue}}
\subsubsectionfont{\color{nalgo}}

%Librías tikz
\usepackage{pgf,tikz}
\usepackage{mathrsfs}
\usetikzlibrary{arrows}
\usetikzlibrary[patterns]
\newcommand{\degre}{\ensuremath{^\circ}}
\definecolor{qqwuqq}{rgb}{0.,0.39215686274509803,0.}
\definecolor{ffttww}{rgb}{1.,0.2,0.4}
%Hipervinculos
\usepackage{hyperref}

\usepackage{fancyhdr}
\pagestyle{fancy}
\fancyhead[L]{Combinaciones}
\fancyhead[C]{Licenciatura en ciencias de la computación}
\fancyhead[R]{USACH}

%interlineado
\renewcommand{\baselinestretch}{1.2}

%\bibitem{Yahoo} \textsc{Andres G} (2009),
%\textbf{¿Generar números aleatorios negativos en Lenguaje C?} En \textsc{Yahoo! respuestas}
%Recuperado el el 23 del julio del 2014
%\url{https://es.answers.yahoo.com/question/index?qid=20091121055249AAUQH3N}

\newcommand{\biblio}[7]{
\bibitem{#1} \textsc{#2} (#3),
\textbf{#4} En \textsc{#5}
Recuperado el #6
\url{#7}
}

% Last, F. M. (Year Published) Book. City, State: Publisher.
\newcommand{\book}[5]{
\bibitem{#1} \textsc{#2} (#3),
\textbf{#4}  \textsc{#5} Estado: Publicado
}

\begin{document}

\begin{titlepage}

\begin{center}

{\Large { Licenciatura en ciencia de la computación} }

\includegraphics[scale=1]{UDSCNRJ}
\\[1cm]

{\Huge \textsc{Combinaciones}}\\[0.7cm]

{\huge  Matemática Computacional}\\[2cm]


\begin{minipage}[l]{0.4\textwidth}
	\begin{flushleft}
	\linespread{1}
		\textbf{\textsf{Profesor:}}\\
		\large Nicolas Thériault
	\end{flushleft}
\end{minipage}
\begin{minipage}[l]{0.4\textwidth}

	\begin{flushright}

		\textbf{\textsf{Autor:}}\\
		\linespread{1}
		\large Sergio Salinas\\
		\large Danilo Abellá\\

	\end{flushright}
\end{minipage}

\end{center}

\end{titlepage}

\section*{Introducción}

Pico pal que lee
\section{Algoritmo Implementado}

\section{Estrategía uno}

\begin{algorithm}[H]
 \KwData{n,r $\in$ $\mathbb{Z}$}
 \KwResult{C(n,r)}
	nr = n -r\;
	n\_result = n! \;
	r\_result = r! \;
	nr\_result = nr! \;
	combinatoria = $\dfrac{n\_result}{r\_result \cdot nr\_result}$ \;
 \caption{How to write algorithms}
\end{algorithm}

\subsection{Estrategía dos}

Para la creación de este algoritmo se baso en el siguiente ejemplo:
\[\begin{matrix}
{7 \choose 4}  &= \dfrac{7!}{4! \cdot 3!} = \dfrac{7 \cdot 6 \cdot 5}{3 \cdot 2}\\ 
 &\\
{7 \choose 3}  &= \dfrac{7!}{3!\cdot4!} = \dfrac{7 \cdot 6 \cdot 5}{3 \cdot 2 }\\
\end{matrix}
\]

Gracias a la propiedad de simetría da el mismo resultado pero con los resultados cambiados, cuando $n-r > r$ ocurre que el mayor número del denominador se queda a la derecha en caso contrario el mayor se queda la izquierda.

Además se eso se observa cuando se factoriza factoriales el númerador queda de la forma $n \cdot (n-1) \cdot \ldots \cdot (r+1)!$, en base a todo esto se creo el siguiente algoritmo

\begin{itemize}
\item Se usa la propiedad se simetría para así lograr que r tenga su máximo número. Ejemplo:

Sí 

\item
\end{itemize}


%\begin{algorithm}[H]
% \KwData{n,r $\in$ $\mathbb{Z}$}
% \KwResult{C(n,r)}
% 	nr = n-r
%	\eIf{n - r $>$ r}{
%   		r=r+1;
%   		$n_result = r \cdot r+1 \cdot \ldots \cdot n$
%   		nr_result = (nr)!
%   		comb = $\frac{n_result}{nr_result}$
%   		}{
%   		nr=nr+1;
%   		$n_result = nr \cdot nr+1 \cdot \ldots \cdot n$
%   		r_result = r!
%   		comb = $\frac{n_result}{r_result}$
%  	}
% \caption{How to write algorithms}
%\end{algorithm}


\subsection{Estrategía tres}

\subsection{Estregía cuatro}

\subsection{Estrategía cinco}



\section{Formulación experimentos}

Para probar los experimentos se hizo un script en bash que compilara una vez y ejecutara varias veces el ejecutable, variando los metodos de entrada. Los experimentos son.

\subsection{n fijo y r varia}

En este experimento la función combinatoria tiene un n con valor 1000 y parte con r de valor 1, se ejecuta el ejecutable varias veces incrementandole valores a r hasta sea igual que n.

\subsection{n varia y r es fijo}

Para este experimento se asigna un valor inicial a n y r de 100, n  va creciendo hasta llegar a 1000.

\subsection{n fijo y r siendo constantemente multiplicado con k}


En este experimento de deja un n fijo de 10000 y un r de 1, se ejecuta varias veces multiplicando constantemente r por una constante k igual 2 hasta que r sea igual o mayor que r.

\subsection{Factorial de (n/r) con r<n}

En este experimento se utilizó la combinación numérica con la estrategia 2.2 mediante la cual sacamos el valor de una combinatoria utilizando el factorial de la división entre la variable n y r respectivamente , con r < n.
Cabe mensionar que esto es una aproximación numérica.


\subsection{Aproximación de la factorial de n,r y (n-r) mediante la fórmula de Striling}

Se utilizó propiedades de potencia, división y multiplicación en este experimento, también cave mencionar que es un resultado apróximado de la factorial.

\section{Información de Hardware y Software}

PC - Danilo Abellá

Software
SO: Ubuntu 16.04.1 LTS
GMP Library
Mousepad 0.4.0

Hardware



\section{Curvas de desempeño de resultados}

\section{Conclusiones}


\end{document}
